\chapter{Organização dos Arquivos}

\section{Entradas públicas}

A entrada da monografia é \texttt{tcc.tex}. Para consultar recursos demonstrativos,
use \texttt{tcc-exemplo-undf.tex} e a pasta \texttt{exemplo/}. Não desenvolva o
trabalho diretamente sobre o exemplo completo, pois ele contém elementos
opcionais e textos destinados a teste.

Na branch de desenvolvimento, as demais entradas são \texttt{tcc-artigo.tex}
(TCC-artigo), \texttt{artigo.tex} (artigo autônomo),
\texttt{relatorio-projeto-aplicado.tex} (relatório parcial de cada projeto),
\texttt{portfolio.tex} (relatório final do portfólio, modalidade 1.1),
\texttt{relatorio-projeto.tex}
(relatório final de projeto técnico/profissional, modalidade 1.3) e
\texttt{negocio-tecnologico.tex} (relatório final de negócio tecnológico,
modalidade 1.4). O índice \texttt{modelos/README.md} relaciona cada modelo ao
seu exemplo e aos arquivos editáveis. A base pública possui desenvolvimento
independente e não inclui automaticamente essas modalidades experimentais.

O relatório parcial preserva memória, conteúdo e evidências de cada projeto.
Sua abertura é enxuta, sem capa de TCC, aprovação, resumos ou sumário; não
exige periodicidade, parecer ou cronograma. Edite a identificação, o relato
e as fontes consultadas em \texttt{modelos/projeto-aplicado/}.

O portfólio final reúne todos os projetos do percurso declarado: introdução,
organização, inventário, síntese de cada projeto, análise integradora e
conclusão. Não basta concatenar os registros. Edite os metadados, resumos,
conteúdo e bibliografia em \texttt{modelos/portfolio/}. O exemplo preenchido
está em \texttt{portfolio-exemplo-undf.tex}, com arquivos em
\texttt{exemplo/portfolio/}. O número de projetos do exemplo é demonstrativo.

Os elementos obrigatórios e opcionais e as referências do portfólio final
preservam o padrão do TCC convencional. Artigo científico não é exigido
nesses modelos nem no relatório profissional. O documento de agosto de 2026
foi referência de organização do relato conclusivo. A admissibilidade e o
recorte do percurso devem ser confirmados com o curso.

\section{Mapa do projeto}

\begin{longtable}{p{0.31\linewidth}p{0.57\linewidth}}
\toprule
Caminho & Finalidade \\
\midrule
\endfirsthead
\toprule
Caminho & Finalidade \\
\midrule
\endhead
\texttt{tcc.tex} & Entrada do modelo em branco e ordem de montagem do documento. \\
\texttt{capitulos/} & Conteúdo textual da monografia. \\
\texttt{obrigatorio/} & Resumos e base bibliográfica da monografia. \\
\texttt{artigo/} & Metadados, resumos e miolo dos modelos de artigo. \\
\texttt{modelos/} & Índice e arquivos preenchíveis de relatório parcial, portfólio, projeto profissional e negócio tecnológico; apresentação comum dos TCCs finais. \\
\texttt{exemplo/} & Conteúdo didático independente dos modelos para preenchimento. \\
\texttt{opcional/} & Errata, dedicatória, agradecimentos, epígrafe, glossário, apêndices e anexos. \\
\texttt{figuras/} & Imagens usadas pelo trabalho. \\
\texttt{identidade/} & Ativos institucionais carregados pelo template. \\
\texttt{src/core/} & Metadados, estrutura, títulos, layout, arquivos e referências cruzadas. \\
\texttt{src/abnt/} & Camadas técnicas organizadas por norma. \\
\texttt{src/identity/} & Capa e páginas de identificação. \\
\texttt{src/biblatex/} & Estilo e comandos bibliográficos. \\
\bottomrule
\end{longtable}

\section{Arquivos que o usuário edita}

Edite \texttt{tcc.tex}, os arquivos de \texttt{capitulos/}, \texttt{obrigatorio/}, \texttt{opcional/} e os ativos inseridos em \texttt{figuras/}. Arquivos internos de \texttt{src/} devem ser modificados apenas durante manutenção do template e acompanhados de testes nos modelos em branco e de exemplo.

Para as demais modalidades, edite a entrada escolhida e a pasta indicada no
índice dos modelos. Mantenha os arquivos de exemplos separados do trabalho real.
No TCC, os dados catalográficos continuam obrigatórios na versão final; o caminho
\texttt{opcional/ficha-catalografica.tex} é mantido por compatibilidade e não
define a obrigatoriedade do elemento.

\section{Inclusão de capítulos}

O arquivo principal usa \texttt{\textbackslash undfincludechapter}. Para ativar a seção opcional de dados ou corpus, retire o comentário da linha correspondente:

\begin{lstlisting}
\undfincludechapter{capitulos/04-dados}
\end{lstlisting}

Para adicionar outra seção primária, crie o arquivo na pasta \texttt{capitulos/} e inclua-o na posição desejada. Evite renomear arquivos internos apenas por preferência estética; caminhos inconsistentes são causa frequente de falha.

\section{Arquivos auxiliares}

LuaLaTeX e Biber criam arquivos como \texttt{.aux}, \texttt{.bcf}, \texttt{.bbl}, \texttt{.log}, \texttt{.toc}, \texttt{.lof} e \texttt{.lot}. Eles não devem ser editados manualmente. Use limpeza controlada quando houver estado inconsistente e recompile o documento completo.
